% Part: first-order-logic
% Chapter: sequent-calculus
% Section: proof-theoretic-notions

\documentclass[../../../include/open-logic-section]{subfiles}

\begin{document}

\begin{editorial}
  Deze paragraaf bundelt de definities van de bewijsbaarheidsrelatie en
  consistentie voor natuurlijke deductie.
\end{editorial}

\iftag{FOL}
      {\olfileid{fol}{seq}{ptn}}
      {\olfileid{pl}{seq}{ptn}}

\olsection{Bewijstheoretische begrippen}

\begin{explain}
Zoals we enkele belangrijke semantische begrippen hebben gedefinieerd
(geldigheid, semantisch gevolg en vervulbaarheid), definiëren we nu de
overeenkomstige \emph{bewijstheoretische begrippen}. Deze worden niet
gedefinieerd aan de hand van de vervulling van !!{sentence}s in
!!{structure}s, maar aan de hand van de !!{derivability} of
!!{nonderivability} van bepaalde sequenten. De ontdekking dat deze
begrippen samenvallen was van groot belang. Dat ze samenvallen, is de
inhoud van de \emph{correctheidsstelling} en de
\emph{volledigheidsstelling}.
\end{explain}

\begin{defn}[Stellingen]
!!^a{sentence}~$!A$ is een \emph{stelling} als er in~$\Log{LK}$
!!a{derivation} van de sequent $\quad \Sequent !A$ bestaat. We schrijven
$\Proves !A$ als $!A$ een stelling is en $\Proves/ !A$ als dat niet zo
is.
\end{defn}

\begin{defn}[!!^{derivability}]
!!^a{sentence}~$!A$ is \emph{!!{derivable} uit} een verzameling
!!{sentence}s~$\Gamma$, genoteerd $\Gamma \Proves !A$, dan en slechts
dan als er een eindige deelverzameling~$\Gamma_0 \subseteq \Gamma$ en
een rij~$\Gamma_0'$ van de !!{sentence}s in~$\Gamma_0$ bestaan zodat
$\Log{LK}$ de sequent $\Gamma_0' \Sequent !A$ !!{derive}s. Als $!A$
niet !!{derivable} is uit $\Gamma$, schrijven we $\Gamma \Proves/ !A$.
\end{defn}

Vanwege de contractie-, verzwakkings- en verwisselingsregels zijn de
volgorde en het aantal voorkomens van de !!{sentence}s in~$\Gamma_0'$
niet van belang: als $\Gamma_0' \Sequent !A$ !!{derivable} is, dan
geldt dit ook voor $\Gamma_0'' \Sequent !A$ voor elke $\Gamma_0''$ die
dezelfde !!{sentence}s bevat als~$\Gamma_0'$. Als bijvoorbeeld
$\Gamma_0 = \{!B, !C\}$, dan zijn zowel $\Gamma_0' = \tuple{!B, !B,
!C}$ als $\Gamma_0'' = \tuple{!C, !C, !B}$ rijen die uitsluitend de
!!{sentence}s uit~$\Gamma_0$ bevatten. Is een sequent met de ene rij
!!{derivable}, dan is de overeenkomstige sequent met de andere rij dat
ook, bijvoorbeeld:
\begin{prooftree}
  \AxiomC{}
  \Deduce$!B, !B, !C \fCenter !A$
  \RightLabel{\LeftR{\Contraction}}
  \UnaryInf$!B, !C \fCenter !A$
  \RightLabel{\LeftR{\Exchange}}
  \UnaryInf$!C, !B \fCenter !A$
  \RightLabel{\LeftR{\Weakening}}
  \UnaryInf$!C, !C, !B \fCenter !A$
\end{prooftree}
Voortaan bedoelen we, voor een eindige verzameling
!!{sentence}s~$\Gamma_0$, met $\Gamma_0 \Sequent !A$ elke sequent waarvan het
antecedent een rij van !!{sentence}s uit~$\Gamma_0$ is. Waar nodig
veronderstellen we stilzwijgend toepassingen van contractie,
verwisseling en verzwakking.

\begin{defn}[Consistentie]
Een verzameling zinnen~$\Gamma$ is \emph{inconsistent} dan en slechts
dan als er een eindige deelverzameling~$\Gamma_0 \subseteq \Gamma$
bestaat waarvoor $\Log{LK}$ de sequent $\Gamma_0 \Sequent \quad$
!!{derive}s. Is $\Gamma$ niet inconsistent, dat wil zeggen: geldt voor
elke eindige $\Gamma_0 \subseteq \Gamma$ dat $\Log{LK}$ de sequent
$\Gamma_0 \Sequent \quad$ niet !!{derive}s, dan heet de verzameling
\emph{consistent}.
\end{defn}

\begin{prop}[Reflexiviteit]
\ollabel{prop:reflexivity}
Als $!A \in \Gamma$, dan $\Gamma \Proves !A$.
\end{prop}

\begin{proof}
De beginsequent $!A \Sequent !A$ is !!{derivable}, en $\{!A\}
\subseteq \Gamma$.
\end{proof}

\begin{prop}[Monotonie]
\ollabel{prop:monotonicity}
Als $\Gamma \subseteq \Delta$ en $\Gamma \Proves !A$, dan $\Delta
\Proves !A$.
\end{prop}

\begin{proof}
Stel $\Gamma \Proves !A$. Er is dus een eindige $\Gamma_0 \subseteq
\Gamma$ waarvoor $\Gamma_0 \Sequent !A$ !!{derivable} is. Omdat
$\Gamma \subseteq \Delta$, is $\Gamma_0$ ook een eindige
deelverzameling van~$\Delta$. De !!{derivation} van $\Gamma_0 \Sequent
!A$ toont daarom ook aan dat $\Delta \Proves !A$.
\end{proof}

\begin{prop}[Transitiviteit]
\ollabel{prop:transitivity}
Als $\Gamma \Proves !A$ en $\{!A\} \cup \Delta \Proves !B$, dan
$\Gamma \cup \Delta \Proves !B$.
\end{prop}

\begin{proof}
Als $\Gamma \Proves !A$, bestaan er een eindige $\Gamma_0 \subseteq
\Gamma$ en !!a{derivation}~$\pi_0$ van $\Gamma_0 \Sequent !A$. Als
$\{!A\} \cup \Delta \Proves !B$, bestaat er voor een eindige
deelverzameling $\Delta_0 \subseteq \Delta$ !!a{derivation}~$\pi_1$
van $!A, \Delta_0 \Sequent !B$. Beschouw de volgende !!{derivation}:
\begin{prooftree}
  \AxiomC{}
  \RightLabel{$\pi_0$}
  \Deduce$\Gamma_0 \fCenter !A$
  \AxiomC{}
  \RightLabel{$\pi_1$}
  \Deduce$!A, \Delta_0 \fCenter !B$
  \RightLabel{\Cut}
  \BinaryInf$\Gamma_0, \Delta_0 \fCenter !B$
\end{prooftree}
Omdat $\Gamma_0 \cup \Delta_0 \subseteq \Gamma \cup \Delta$, toont
dit aan dat $\Gamma \cup \Delta \Proves !B$.
\end{proof}

Dit betekent in het bijzonder dat uit $\Gamma \Proves !A$ en $!A
\Proves !B$ volgt dat $\Gamma \Proves !B$. Ook volgt dat, als $!A_1,
\dots, !A_n \Proves !B$ en $\Gamma \Proves !A_i$ voor elke~$i$, dan
$\Gamma \Proves !B$.

\begin{prop}
\ollabel{prop:incons}
$\Gamma$ is inconsistent dan en slechts dan als $\Gamma \Proves {!A}$
voor elke zin~$!A$.
\end{prop}

\begin{proof}
Opgave.
\end{proof}

\begin{prob}
Bewijs \olref[fol][seq][ptn]{prop:incons}.
\end{prob}

\begin{prop}[Compactheid]
  \ollabel{prop:proves-compact}
  \begin{enumerate}
  \item Als $\Gamma \Proves !A$, dan is er een eindige deelverzameling
    $\Gamma_0 \subseteq \Gamma$ waarvoor $\Gamma_0 \Proves !A$.
  \item Als elke eindige deelverzameling van~$\Gamma$ consistent is,
    dan is $\Gamma$ consistent.
  \end{enumerate}
\end{prop}

\begin{proof}
  \begin{enumerate}
    \item Als $\Gamma \Proves !A$, dan is er een eindige
      deelverzameling $\Gamma_0 \subseteq \Gamma$ waarvoor de sequent
      $\Gamma_0 \Sequent !A$ !!a{derivation} heeft. Bijgevolg geldt
      $\Gamma_0 \Proves !A$.
    \item Als $\Gamma$ inconsistent is, dan is er een eindige
      deelverzameling~$\Gamma_0 \subseteq \Gamma$ waarvoor $\Log{LK}$
      de sequent $\Gamma_0 \Sequent \quad$ !!{derive}s. Dan is
      $\Gamma_0$ dus een inconsistente eindige deelverzameling
      van~$\Gamma$.
  \end{enumerate}
\end{proof}
\end{document}
